% Section 8: The pseudoscalar} | Part: 1 | Depends on: see docs/REWRITE_PLAN.md | Status: todo | Original lines: 1536-1554
%!TEX root = ../../../main.tex
\subsection{The pseudoscalar}

Let $(e_1,\dots,e_n)$ be an orthonormal basis of a non-degenerate space $V=\R^{p,q}$.

\begin{definition}[Pseudoscalar]
The \emph{unit pseudoscalar} is $\pss=e_1e_2\cdots e_n=e_1\wedge\cdots\wedge e_n\in\GA^n(V)$. Changing the basis by an orientation-preserving orthogonal map leaves $\pss$ unchanged; an orientation-reversing one changes its sign. Choosing $\pss$ is therefore the same as choosing an \emph{orientation} of $V$.
\end{definition}

\begin{proposition}\label{prop:pss}
\begin{enumerate}
  \item $\pss^2=(-1)^{n(n-1)/2}(-1)^q$ and $\pss^{-1}=\rev\pss/(\pss\rev\pss)$. In the Euclidean case $\pss^{-1}=\rev\pss=(-1)^{n(n-1)/2}\pss$.
  \item $e_i\pss=(-1)^{n-1}\pss e_i$: $\pss$ commutes with all vectors if $n$ is odd, and anticommutes with them if $n$ is even. In particular $\pss$ is central in $\GA(V)$ when $n$ is odd.
  \item For any $k$-vector $A$ the product $A\pss$ is a pure $(n-k)$-vector.
\end{enumerate}
\end{proposition}
\begin{proof}
(1) $\rev\pss=e_n\cdots e_1=(-1)^{n(n-1)/2}\pss$ and $\pss\rev\pss=\prod_i e_i^2=(-1)^q$, so $\pss^2=(-1)^{n(n-1)/2}\pss\rev\pss$. (2) Moving $e_i$ through $\pss$ it anticommutes with the $n-1$ factors $e_j$, $j\neq i$, and commutes with $e_i$. (3) For a basis blade $e_J$ with $|J|=k$, $e_J\pss$ equals, up to a sign and a product of $\varepsilon_j$, the basis blade $e_{J^c}$ of the complementary indices.
\end{proof}
